% Part: lambda-calculus
% Chapter: lambda-definability
% Section: pairs

\documentclass[../../../include/open-logic-section]{subfiles}

\begin{document}

\olfileid{lam}{ldf}{pai}
\olsection{زوج‌ها و تابع پیشین}

\begin{defn}
زوج~$M$ و~$N$ (با نماد~$\tuple{M,N}$) چنین تعریف می‌شود:
\[
\tuple{M,N} \ident \lambd[f][fMN].
\]
\end{defn}
  
به‌طور شهودی، زوج تابعی است که تابعی را می‌پذیرد و آن تابع را بر دو
مؤلفهٔ زوج اعمال می‌کند. بر پایهٔ این ایده، سازندهٔ زیر را داریم که دو
ترم را می‌پذیرد و زوجی شامل آن‌ها بازمی‌گرداند:
\[
  \fn{Pair} \ident \lambd[mn][\lambd[f][fmn]]
\]
همچنین می‌خواهیم با داشتن یک زوج بتوانیم مؤلفه‌های آن را بازیابی کنیم.
برای این کار به دو تابع دسترسی نیاز داریم که یک زوج را به‌عنوان آرگومان
می‌پذیرند و مؤلفهٔ نخست یا دوم آن را بازمی‌گردانند:
\begin{align*}
  \fn{Fst} & \ident \lambd[p][p(\lambd[mn][m])]\\
  \fn{Snd} & \ident \lambd[p][p(\lambd[mn][n])]
\end{align*}

\begin{prob}
  توضیح دهید چرا توابع دسترسی~$\fn{Fst}$ و~$\fn{Snd}$ درست کار می‌کنند.
\end{prob}

اکنون با استفاده از زوج‌ها می‌توانیم تابع پیشین را~!!{lambda define}:
\[
  \fn{Pred} \ident \lambd[n][\fn{Fst}(n (\lambd[p][\tuple{\fn{Snd}\, {p}, \fn{Succ}(\fn{Snd}\, {p})}]) \tuple{\num 0, \num 0})]
\]
به یاد آورید که~$\num n\, f x$ به~$f^{n}(x)$ کاهش می‌یابد. در اینجا
$f$ تابعی است که زوج~$p$ را می‌پذیرد و زوج تازه‌ای شامل مؤلفهٔ دوم~$p$
و جانشین مؤلفهٔ دوم را بازمی‌گرداند؛ و~$x$ زوج
$\tuple{\num 0,\num 0}$ است. بنابراین برای~$n=0$ حاصل
$\tuple{\num 0,\num 0}$، و در غیر این صورت
$\tuple{\num{n-1}, \num n}$ است. سپس~$\fn{Pred}$ مؤلفهٔ نخست حاصل را
بازمی‌گرداند.

تفریق را می‌توان با اعمال~$b$بارهٔ~$\fn{Pred}$ بر~$a$ تعریف کرد:
\[
\fn{Sub} \ident \lambd[ab][b \fn{Pred}\, a].
\]
    
\end{document}
